Using these little Lemmas, we can state the main result of this chapter. Roughly speaking, it says the following: every initial value problem that is locally Lipschitz has a unique solution. This solution can either be extended to all of $\mathbb{R}$ as a domain or until it explodes so that we couldn't define a solution outside of our solution interval.

\begin{remark}[Locally Lipschitz]
	It should be clear that any continuously differentiable vector field $F \in C^{1}(U, \mathbb{R}^{n})$ is always locally Lipschitz, i.e. around every point $p \in U$, there exists a ball $B_{r}(p) \subset U$ such that $F|_{B_{r}(p)}$ is Lipschitz. Indeed, since $DF(p)$ varies continuously, we can restrict $F$ to a closed ball $\overline{ B_{2r}(p) }$, where hence $\lVert DF(p) \rVert_{2}$ attains its extrema. But then we can just apply \ref{prop:derivative-lipschitz} (modify for vector valued $F$ as an exercise).
\end{remark}
% ignore label
\begin{prop}[Local existence and uniqueness via Picard]\label{prop:peano-ex} 
	Let $r > 0$, $(t_{0}, x_{0}) \in \mathbb{R} \times \mathbb{R}^{n}$ and consider
	\begin{align*}
		F: (t_{0} - r, t_{0}+r) \times B_{r}(x_{0}) \to \mathbb{R}^{n}
	\end{align*}
	continuous. Suppose $\exists C \geq 1, L > 0$ such that
	\begin{align*}
		\left| F(t,x)  \right| \leq C, \quad \left| F(t, x_{1}) - F(t, x_{2}) \right| \leq L \left| x_{1}-x_{2} \right|
	\end{align*}
	for all $t \in (t_{0}-r, t_{0}+r)$ and $x, x_{1},x_{2} \in B_{r}(x_{0})$.

	Then, for any $\delta > 0$ with $\delta < \min\left\{ \tfrac{r}{2C}, \tfrac{1}{2L}\right\}$, there exists a \textit{unique} differentiable function $u: (t_{0} - \delta, t_{0} + \delta) \to B_{r}(x_{0})$ satisfying
	\begin{align*}
		u(t_{0}) = x_{0}, \quad u'(t) = F(t,u(t)).
	\end{align*}
\end{prop}
\begin{proof}
	Choose $\delta > 0$ as proposed and denote $I = (t_{0} - \delta, t_{0}+\delta)$. We can rewrite the differential equation as
	\begin{align*}
		u(t) = x_{0} + \int_{t_{0}}^{t} F(s,u(s))  \, ds,
	\end{align*}
	where we integrate component-wise. If we could make this into a fixed point equation for a fitting Lipschitz contraction, we'd get the result from Banach Fixed Point Theorem \ref{thm:banach_fixed_point}.

	For that, we define
	\begin{align*}
		V = \left\{ u \in C(I, \mathbb{R}^{n}) \mid \lVert u - x_{0} \rVert_{\infty} \leq r_{0}  \right\},
	\end{align*}
	which is a subset of $C(I, \mathbb{R}^{n})$, which is a complete metric space with supremum norm $\lVert \cdot \rVert_{\infty}$ by the same proof as $C_{b}(I, \mathbb{R}^{n})$. But $V$ is a closed subset of a complete space, is hence itself complete.

	We define the \textit{Picard map} $T: V \to C(I, \mathbb{R}^{n})$ by
	\begin{align*}
		(Tu)(t) = x_{0} + \int_{t_{0}}^{t} F(s,u(s))  \, ds,
	\end{align*}
	and look for the fixed point. First, we have to verify that $T$ maps into $V$. For that, let $u \in V$ and $t \in I$
	\begin{align*}
		\left| Tu(t) - x_{0} \right| = \left| \int_{t_{0}}^{t} F(s,u(s))  \, ds \right| \leq \int_{t_{0}}^{t} \left| F(s, u(s)) \right|  \, ds \leq C \left| t-t_{0} \right| \leq C \delta \leq \frac{r}{2}.
	\end{align*}
	So the Picard map can be restricted to $T: V \to V$. Now we show that it is a Lipschitz contraction by computing
	\begin{align*}
		\left|T u_1(t)-T u_2(t)\right| & =\left|\int_{t_0}^t F\left(s, u_1(s)\right)-F\left(s, u_2(s)\right) d s\right|                                \\
		                               & \leq \int_{t_0}^t\left|F\left(s, u_1(s)\right)-F\left(s, u_2(s)\right)\right| d s                             \\
		                               & \leq \int_{t_0}^t L\left|u_1(s)-u_2(s)\right| d s                                                             \\
		                               & \leq L\left|t-t_0\right| \sup _{s \in\left(t_0-\delta, t_0+\delta\right)}\left|\left(u_1-u_2\right)(s)\right| \\
		                               & \leq L \delta\left\|u_1-u_2\right\|<\frac{1}{2}\left\|u_1-u_2\right\|
	\end{align*}
    for all $t \in I$. Hence $T$ is Lipschitz with constant $\frac{1}{2}$ and Banach applies. Note that Banach also gives us uniqueness.
\end{proof}

\begin{theorem}[Picard-Lindelöf (Cauchy-Lipschitz)]\label{thm:picard}
	Let $U \subset \mathbb{R} \times \mathbb{R}^{n}$ open, $(t_{0}, x_{0}) \in U$ and $F: U \to \mathbb{R}^{n}$ be continuous.

	Assume that $F$ is "locally Lipschitz in space". That is, for every $(t_{0}, x_{0}) \in U$, $\exists \epsilon > 0$ and $L \geq 0$ such that
	\begin{align*}
		\left| F(t,x) - F(t,x') \right| \leq L \left| x-x' \right| \quad \forall (t,x), (t,x') \in B_{\epsilon}((t_{0}, x_{0})) \cap U.
	\end{align*}

	Then, $\exists$ an interval $I = (a,b) \subset \mathbb{R}$ with $t_{0} \in I$ and a differentiable function $u: I \to \mathbb{R}^{d}$ such that the following holds:
	\begin{enumerate}
		\item $(t,u(t)) \in U$ for all $t \in I$ and $u$ solves the initial value problem
		      \begin{align*}
			      u'(t) = F(t,u(t)), \quad u(t_{0}) = x_{0}.
		      \end{align*}
		\item For any other solution $v: J \to \mathbb{R}^{n}$ defined on an open interval $J$ of the same initial value problem with $t_{0} \in J$, we have $J \subseteq I$ and $u|_{J} = v$.
		\item The limits $\lim_{t \to a}(t,u(t))$ and $\lim_{t \to b}(t,u(t))$ either do not exist or are not in $U$.
	\end{enumerate}
\end{theorem}
\begin{proof}
	Before we begin, one convinces themselves that locally, $F$ satisfies the assumptions in \ref{prop:peano-ex}, hence we can apply the proposition repeatedly. I won't be writing out all the qualifiers each time I invoke it.

    First we tackle uniqueness. Let $u_{1}: I_{1} \to \mathbb{R}^{n}$ and $u_{2}: I_{2} \to \mathbb{R}^{n}$ satisfy
    \begin{align*}
        u_{i}'(t) =F(t, u_{i}(t)), \quad u_{i}(t_{0}) = x_{0}
    \end{align*}
    on open intervals $I_{1}, I_{2} \subset \mathbb{R}$ with $t_{0} \in I = I_{1} \cap I_{2}$ and claim $u_{1}(t) = u_{2}(t)$ on $I$. 

    For that, consider
    \begin{align*}
        S = \left\{ t \in I \mid u_{1}(t) = u_{2}(t) \right\}.
    \end{align*}
    The complement $I \setminus S$ is open by continuity of $u_{1}, u_{2}$ and $S$ is not empty because $t_{0} \in S$. For any $t_{1} \in S$, both $u_{1}$ and $u_{2}$ solve the initial value problem with initial conditions $u_{1}(t_{1}) = u_{2}(t_{1}) =: x_{1}$. But Proposition \ref{prop:peano-ex} guarantees an open set around $t_{1}$ where $u_{1}$ and $u_{2}$ align. Hence $S$ is open. Since $I$ is an interval, it is connected and $I \setminus S = \emptyset$.

    Now we turn to existence. Consider $\mathcal{J}$ the set of pairs $(J, v)$, where $J \subset \mathbb{R}$ is an open interval with $t_{0} \in J$ and a function $u: J \to \mathbb{R}^{n}$ that solves the initial value problem. By \ref{prop:peano-ex}, $\mathcal{J}$ is non-empty. We set $I = \bigcup_{(J,v) \in \mathcal{J}} J$ and define $u: T \to \mathbb{R}^{n}$ piecewise, which is well-defined by uniqueness. $I$ is also open as a union of open sets and solves the initial problem.

    It remains to show the behavior at the end points $a$ and $b$. We show $b$, $a$ is analogous. If $b = \infty$, we are done. Otherwise, let $b = \sup I < \infty$ and assume 
    \begin{align*}
        \lim_{t \to b}(t,u(t)) = (b,x_{b}) \in U
    \end{align*}
    exists. Then, by \ref{prop:peano-ex}, we have a solution $w$ to $w(b) = x_{b}$ and $w'(t) = F(t,w(t))$ defined on $J = (b - \delta, b + \delta)$ for some $\delta > 0$. 

    We use this to extend $I$, leading to a contradiction. We define $v: (a,b+\delta) \to \mathbb{R}^{n}$ by
    \begin{align*}
        v(t) = \begin{cases}
            u(t) \quad &\text{if } t \in (a,b) \\ 
            w(t) \quad &\text{if }t \in [b, b+\delta)
        \end{cases}
    \end{align*}
    $v$ is continuous by construction, also at $b$. We know $v'(t) = F(t,v(t))$ holds on all $(a, b+\delta)$ maybe except at $t=b$. 

    It remains to show $v$ is differentiable at $b$ and $v'(b) = F(b,v(b))$. From the right,
    \begin{align*}
        \lim_{h \to 0^{+}} \frac{v(b+h) - v(b)}{h} = w'(b) = F(b,w(b)) = F(b, v(b)). 
    \end{align*}
    From the left, using component-wise L'Hopital,
    \begin{align*}
        \lim_{h \to 0^{-}} \frac{u(b+h) - u(b) - F(b, u(b)) \cdot h}{h} = \lim_{h \to 0^{-}} u'(b+h) - F(b, u(b)) = 0,
    \end{align*}
    hence $u'(b) = F(b,u(b))$. But this proves the theorem.
\end{proof}
